%% Question bank chapter file - generated by bank_tools.py build (do not edit while the bank is still being built:
%% it is regenerated from the per-test data in Question Bank\data). Every question is one bankq block:
%%   \begin{bankq}[id=..., source={...}, paper=.., qno=.., type=..., subtopic={...}, answer={...}, answer_src=key|solved,
%%                 working={...}, difficulty=1-10, diff_parts={K S I C F}, diff_note={...},
%%                 group=<passage id or empty>, also_in={duplicate copies in other tests}]
%% difficulty = min(10, K+S+I+C+F-1): K concepts 1-3, S steps 1-4, I insight 0-2, C computation 0-2, F format 0-1
%%   ...question LaTeX, options, figures (\bfig / \bfigi -> img/)...
%%   \end{bankq}
%% Passages shared by several questions are bankpassage blocks; their questions carry group=<passage id>.
\documentclass[10pt]{article}
\usepackage[a4paper,margin=1.8cm]{geometry}
\usepackage[T1]{fontenc}
\usepackage{amsmath,amssymb,mathtools}
\usepackage{graphicx}
\usepackage[export]{adjustbox}
\usepackage{xcolor,needspace,keyval}
\usepackage{enumitem}
\usepackage{booktabs,array,tabularx}
\usepackage[hidelinks]{hyperref}
\setlength{\parindent}{0pt}
\setlength{\parskip}{3pt}
\makeatletter
\define@key{bq}{id}{\def\bq@id{#1}}\define@key{bq}{source}{\def\bq@src{#1}}\define@key{bq}{paper}{\def\bq@paper{#1}}
\define@key{bq}{qno}{\def\bq@qno{#1}}\define@key{bq}{type}{\def\bq@type{#1}}\define@key{bq}{subtopic}{\def\bq@sub{#1}}
\define@key{bq}{answer}{\def\bq@ans{#1}}\define@key{bq}{answer_src}{\def\bq@asrc{#1}}\define@key{bq}{working}{\def\bq@work{#1}}
\define@key{bq}{group}{\def\bq@group{#1}}\define@key{bq}{also_in}{\def\bq@also{#1}}
\define@key{bq}{difficulty}{\def\bq@diff{#1}}\define@key{bq}{diff_parts}{\def\bq@dparts{#1}}\define@key{bq}{diff_note}{\def\bq@dnote{#1}}
\newenvironment{bankq}[1][]{\def\bq@work{}\def\bq@also{}\def\bq@group{}\def\bq@diff{?}\def\bq@dparts{}\def\bq@dnote{}\setkeys{bq}{#1}%
  \par\needspace{4\baselineskip}\medskip\noindent{\small\textbf{\bq@id}\hfill\color{gray}[\bq@src, Paper \bq@paper, Q\bq@qno]
  \bq@type\ $\cdot$ \bq@sub\ $\cdot$ \textbf{D\bq@diff}}\par\nopagebreak\relax}
  {\par{\small\color{teal}Answer: \textbf{\bq@ans}\ifx\bq@work\@empty\else\ (solved: \bq@work)\fi
  \ifx\bq@also\@empty\else\par\color{gray}Also in: \bq@also\fi}\par}
\define@key{bp}{id}{\def\bp@id{#1}}\define@key{bp}{source}{\def\bp@src{#1}}\define@key{bp}{questions}{\def\bp@qs{#1}}
\newenvironment{bankpassage}[1][]{\setkeys{bp}{#1}\par\needspace{4\baselineskip}\bigskip\noindent
  {\small\textbf{Passage \bp@id}\hfill\color{gray}[\bp@src; questions \bp@qs]}\par\nopagebreak\relax\itshape}{\par}
\makeatother
\newcommand{\bfig}[2]{\par\smallskip{\centering\includegraphics[width=#1pt,max width=0.62\linewidth,max height=0.32\textheight]{img/#2}\par}\smallskip}
\newcommand{\bfigi}[2]{\includegraphics[width=#1pt,max width=\linewidth,max height=0.32\textheight,valign=c]{img/#2}}
\newcommand{\bmissing}[1]{\textit{[figure missing in the source conversion: \detokenize{#1}]}}
\newenvironment{qoptions}{\begin{itemize}[leftmargin=1.8em,itemsep=1pt,topsep=2pt,label={}]}{\end{itemize}}
\newcommand{\opt}[2]{\item[(#1)] #2}
\providecommand{\fwstart}{}\providecommand{\fwend}{}\providecommand{\qfullsplit}{}
\newcolumntype{L}[1]{>{\raggedright\arraybackslash}p{#1}}
\newcolumntype{C}[1]{>{\centering\arraybackslash}p{#1}}
\newcolumntype{R}[1]{>{\raggedleft\arraybackslash}p{#1}}
\providecommand{\sectionhdr}[1]{}\providecommand{\subsectionhdr}[1]{}


\begin{document}
{\Large\bfseries BT --- Question Bank}\par\smallskip
{\small Mathematics chapter M05. Questions from all converted test papers, grouped by source series. Answers: printed key where it matches the paper, otherwise solved (answer\_src=solved, with working). Exact repeats across papers appear once, with also\_in listing the other copies.}\par\medskip\hrule

\section*{TOAS papers}
\begin{bankq}[id=TOAS01-P1-MATH-S1-Q04, source={TOAS - Test 01 (30 Jan 2022)}, paper=1, qno=4, type=single, subtopic={Multinomial}, answer={D}, answer_src=key, difficulty=5, diff_parts={K1 S3 I1 C1 F0}, diff_note={Multinomial cases with 2i+j=l, 3i+2j+k=5 summed}]
The term independent of $x$ in the expansion of $\left(x^2+x+1+\dfrac1x\right)^5$ is equal to -
\begin{qoptions}
\opt{A}{98}
\opt{B}{99}
\opt{C}{100}
\opt{D}{101}
\end{qoptions}
\end{bankq}
\begin{bankq}[id=TOAS01-P1-MATH-S2-Q05, source={TOAS - Test 01 (30 Jan 2022)}, paper=1, qno=5, type=multi, subtopic={General term}, answer={B,C}, answer_src=key, difficulty=5, diff_parts={K2 S3 I0 C0 F1}, diff_note={9th term gives (log2 x)^8=256 so x=4 or 1/4}]
If the $9^{th}$ term in the expansion of $\left(\dfrac{1}{x^2}+\dfrac{x}{2}\log_2x\right)^{12}$ is equal to 495, then the value of $x$ can be-
\begin{qoptions}
\opt{A}{1}
\opt{B}{An integer greater than 1}
\opt{C}{a fraction}
\opt{D}{an irrational number}
\end{qoptions}
\end{bankq}
\begin{bankq}[id=TOAS01-P1-MATH-S2-Q07, source={TOAS - Test 01 (30 Jan 2022)}, paper=1, qno=7, type=multi, subtopic={Coefficient problems}, answer={A,C,D}, answer_src=key, difficulty=6, diff_parts={K1 S3 I1 C1 F1}, diff_note={r C(n,r)=n C(n-1,r-1) then Vandermonde; option A rewritten as 50 C(200,100)}]
Sum of the series $({}^{100}C_1)^2+2.({}^{100}C_2)^2+3.({}^{100}C_3)^2+\dots+100({}^{100}C_{100})^2$ equals-
\begin{qoptions}
\opt{A}{$\dfrac{2^{99}\,[1.3.5.....(199)]}{99!}$}
\opt{B}{$100.{}^{200}C_{100}$}
\opt{C}{$50.{}^{200}C_{100}$}
\opt{D}{$100.{}^{199}C_{99}$}
\end{qoptions}
\end{bankq}
\begin{bankpassage}[id=TOAS01-P1-MATH-S3-G5-6, source={TOAS - Test 01 (30 Jan 2022)}, questions={5-6}]
Let $(1-x+x^2)^{30}=a_0+a_1x+a_2x^2+a_3x^3+\dots\dots\dots+a_{60}x^{60}$.
\end{bankpassage}
\begin{bankq}[id=TOAS01-P1-MATH-S3-Q05, source={TOAS - Test 01 (30 Jan 2022)}, paper=1, qno=5, type=numerical, subtopic={Coefficient problems}, answer={1}, answer_src=key, difficulty=7, diff_parts={K2 S3 I2 C0 F1}, diff_note={Multiply P(x) by P(-x)=(1+x^2+x^4)^30 and read the x^60 coefficient with palindromic symmetry}, group=TOAS01-P1-MATH-S3-G5-6]
$\left(\dfrac{a_0^2-a_1^2+a_2^2-a_3^2+\dots\dots\dots}{a_{30}}\right)$
\end{bankq}
\begin{bankq}[id=TOAS01-P1-MATH-S3-Q06, source={TOAS - Test 01 (30 Jan 2022)}, paper=1, qno=6, type=numerical, subtopic={Coefficient problems}, answer={2}, answer_src=key, difficulty=4, diff_parts={K1 S2 I1 C0 F1}, diff_note={Even-index sum = (P(1)+P(-1))/2}, group=TOAS01-P1-MATH-S3-G5-6]
$4\left(\dfrac{a_0+a_2+a_4+\dots\dots\dots}{(3^{30}+1)}\right)$
\end{bankq}
\begin{bankq}[id=TOAS01-P2-MATH-S1-Q06, source={TOAS - Test 01 (30 Jan 2022)}, paper=2, qno=6, type=single, subtopic={General term}, answer={C}, answer_src=key, difficulty=5, diff_parts={K1 S3 I1 C1 F0}, diff_note={Combine to (x+1/x)(x^2-x^-2)^9 and pick the x^5 term}]
Coefficient of $x^5$ in the expansion of
\[\left(x+\dfrac1x\right)^{10}\left(x-\dfrac1x\right)^{9}\]
is -
\begin{qoptions}
\opt{A}{${}^9C_3$}
\opt{B}{${}^9C_4$}
\opt{C}{$-{}^9C_3$}
\opt{D}{$-{}^9C_4$}
\end{qoptions}
\end{bankq}
\begin{bankq}[id=TOAS01-P2-MATH-S1-Q07, source={TOAS - Test 01 (30 Jan 2022)}, paper=2, qno=7, type=single, subtopic={Coefficient problems}, answer={A}, answer_src=key, difficulty=5, diff_parts={K2 S3 I1 C0 F0}, diff_note={Last 30 coefficients sum to 2^58, then divisors of 58 of form 4k+2}]
If sum of the last 30 coefficients in the expansion of $(1+x)^{59}$, when expanded in ascending power of `$x$' is $2^n$ then number of divisors of `$n$' of the form `$4\lambda+2$' ($\lambda\in N$) is -
\begin{qoptions}
\opt{A}{1}
\opt{B}{0}
\opt{C}{2}
\opt{D}{4}
\end{qoptions}
\end{bankq}
\begin{bankq}[id=TOAS01-P2-MATH-S1-Q08, source={TOAS - Test 01 (30 Jan 2022)}, paper=2, qno=8, type=single, subtopic={Coefficient problems}, answer={B}, answer_src=key, difficulty=2, diff_parts={K1 S2 I0 C0 F0}, diff_note={a1=n and a2=C(n,2)-1 substituted}]
If $(1-x^2+x^3)(1+x)^n=a_0+a_1x+a_2x^2+\dots$, $a_i$'s $\in R$, then the value of $a_1^2-a_1-2a_2$ is -
\begin{qoptions}
\opt{A}{$-2$}
\opt{B}{2}
\opt{C}{4}
\opt{D}{$-4$}
\end{qoptions}
\end{bankq}
\begin{bankq}[id=TOAS01-P2-MATH-S1-Q09, source={TOAS - Test 01 (30 Jan 2022)}, paper=2, qno=9, type=single, subtopic={General term}, answer={C}, answer_src=key, difficulty=2, diff_parts={K1 S2 I0 C0 F0}, diff_note={General term exponent 60-7k=-17}]
The value of `$r$' for which $x^{-17}$ is contained by the $r^{th}$ term in the expansion of $(x^4-x^{-3})^{15}$, is -
\begin{qoptions}
\opt{A}{10}
\opt{B}{11}
\opt{C}{12}
\opt{D}{13}
\end{qoptions}
\end{bankq}
\begin{bankq}[id=TOAS01-P2-MATH-S1-Q10, source={TOAS - Test 01 (30 Jan 2022)}, paper=2, qno=10, type=single, subtopic={Coefficient problems}, answer={D}, answer_src=key, difficulty=3, diff_parts={K1 S2 I1 C0 F0}, diff_note={r C(n,r)=n C(n-1,r-1) then Vandermonde}]
$\displaystyle\sum_{r=0}^{n}r^2\left({}^nC_r\right)^2$ is equal to -
\begin{qoptions}
\opt{A}{$n^2.{}^{2n}C_n$}
\opt{B}{$n^2.{}^{2n}C_{n-1}$}
\opt{C}{$n^2.{}^{2n-2}C_n$}
\opt{D}{$n^2.{}^{2n-2}C_{n-1}$}
\end{qoptions}
\end{bankq}
\begin{bankq}[id=TOAS01-P2-MATH-S1-Q11, source={TOAS - Test 01 (30 Jan 2022)}, paper=2, qno=11, type=single, subtopic={Coefficient problems}, answer={A}, answer_src=key, difficulty=2, diff_parts={K1 S2 I0 C0 F0}, diff_note={Two ratio equations give n=14, r=4}]
If $\dfrac{{}^nC_r}{{}^nC_{(r+1)}}=\dfrac12$ and $\dfrac{{}^nC_{(r+1)}}{{}^nC_{(r+2)}}=\dfrac23$, then -
\begin{qoptions}
\opt{A}{$\dfrac nr=\dfrac72$}
\opt{B}{$\dfrac nr=7$}
\opt{C}{$\dfrac nr=\dfrac73$}
\opt{D}{$\dfrac nr=14$}
\end{qoptions}
\end{bankq}
\begin{bankq}[id=TOAS01-P2-MATH-S2-Q01, source={TOAS - Test 01 (30 Jan 2022)}, paper=2, qno=1, type=numerical, subtopic={Coefficient problems}, answer={7}, answer_src=solved, working={$S(n)=\sum_k(k+2)\binom nk\binom n{k+2}$; by symmetry $k\leftrightarrow n-k-2$ average of $(k+2)$ is $(n+2)/2$ and $\sum\binom nk\binom n{k+2}=\binom{2n}{n-2}$; $S(8)=5\binom{16}6=40040$; $40040/5720=7$}, difficulty=6, diff_parts={K1 S3 I1 C1 F1}, diff_note={k C(n,k)=n C(n-1,k-1) and Vandermonde give S(n)=n C(2n-1,n-2)}]
If $2C_0C_2+3C_1C_3+4C_2C_4+\dots+nC_{n-2}C_n=S(n)$ (where $C_r={}^nC_r$), then the value of $\dfrac{S(8)}{(220)(26)}$ is
\end{bankq}
\begin{bankq}[id=TOAS01-P2-MATH-S2-Q02, source={TOAS - Test 01 (30 Jan 2022)}, paper=2, qno=2, type=numerical, subtopic={Coefficient problems}, answer={5}, answer_src=solved, working={Sum $=\sum_{r\ \text{even}}(r/2+1)\binom{20}r^2=6\cdot\frac12\left[\binom{40}{20}+\binom{20}{10}\right]=3\left[\binom{40}{20}+\binom{20}{10}\right]$, so $8a=40$, $a=5$}, difficulty=8, diff_parts={K2 S3 I2 C1 F1}, diff_note={Even-index sum via alternating sum and symmetry r to 20-r}]
If $1.\left({}^{20}C_0\right)^2+2.\left({}^{20}C_2\right)^2+3.\left({}^{20}C_4\right)^2+\dots+11.\left({}^{20}C_{20}\right)^2$ is equal to $3\left[{}^{8a}C_{4a}+{}^{4a}C_{2a}\right]$, then $a$ is equal to
\end{bankq}
\begin{bankq}[id=TOAS01-P2-MATH-S2-Q03, source={TOAS - Test 01 (30 Jan 2022)}, paper=2, qno=3, type=numerical, subtopic={Coefficient problems}, answer={4}, answer_src=solved, working={$\sum r\binom nr=n2^{n-1}$, so $S=\sum n/2^{n-1}=1/(1-\tfrac12)^2=4$}, difficulty=4, diff_parts={K2 S2 I0 C0 F1}, diff_note={sum r C(n,r)=n 2^(n-1) then arithmetico-geometric series}]
Let $S=\displaystyle\sum_{n=1}^{\infty}\left(\dfrac{n^2}{\sum\limits_{r=0}^{n}r\left({}^nC_r\right)}\right)$, then the value of $S$ is
\end{bankq}
\end{document}
